%===============================================================================
%
%===============================================================================
\section{Les outils de synchronisation}

Allez donc voir votre cours de système d'exploitation centralisé !

%===============================================================================
%      Les outils de synchronisation
%===============================================================================
\section{L'implantation dans ManuX}

   Les moyens de synchronisation entre tâches actuellement disponibles
sont 

\begin{itemize}
   \item les opérations atomiques ;
   \item les exclusions mutuelles.
\end{itemize}

%-------------------------------------------------------------------------------
%      Opérations atomiques
%-------------------------------------------------------------------------------
\subsection{Opérations atomiques}

   Les opérations élémentaires de synchronisation sont définies dans
le fichier {\tt atomique.h}. Elles se basent sur le type
\lstinline!Atomique! et les opérations qui lui sont associées. A part
l'initialisation et la consultation, l'opération importante est la
suivante

\begin{lstlisting}{}
booleen atomiqueTestInit(volatile Atomique * atom, uint32 val, uint32 cond);
\end{lstlisting}

   Elle compare la valeur de l'\lstinline!Atomique! avec
\lstinline!cond! et, en cas d'égalité, affecte l'\lstinline!Atomique!
avec la valeur \lstinline!val! et renvoie \lstinline!TRUE!, sinon elle 
ne fait rien et renvoie \lstinline!FALSE!. Ce qui est important, bien
sur, c'est que la comparaison et l'éventuelle affectation soient
réalisées de façon {\em atomique}.

%-------------------------------------------------------------------------------
%      Exclusions mutuelles
%-------------------------------------------------------------------------------
\subsection{Exclusions mutuelles}

   Les variables d'exclusion mutuelle permettent d'assurer que l'accés
à une ressource (typiquement une section de code critique) se fait de
façon exclusive entre les tâches.

   Le fichier {\tt atomique.h} définit le type et les fonctions
suivants 

\begin{lstlisting}{}
typedef struct _ExclusionMutuelle {
   Atomique   verrou;
   ListeTache tachesEnAttente;
} ExclusionMutuelle;

void initialiserExclusionMutuelle(ExclusionMutuelle * em);
void entrerExclusionMutuelle(ExclusionMutuelle * em);
void sortirExclusionMutuelle(ExclusionMutuelle * em);
\end{lstlisting}

   Naturellement, comme leurs noms l'indiquent, les trois
sous-programmes permettent respectivement d'initialiser une variable
d'exclusion mutuelle, de prendre une exclusion mutuelle et de libérer
une exclusion mutuelle.

